% SGA 5 原 PDF 第 1--12 页；本块不设独立封面及全书目录。
\SGAfrontchapter{原书题名与出版信息}
\SourcePage{1}
\textbf{《数学讲义》（Lecture Notes in Mathematics）589}\par
主编：A.~Dold、B.~Eckmann。\par
\textbf{布瓦玛丽\Term{algebraic-geometry}{代数几何}讨论班，1965—1966（SGA 5）}\par
A.~Grothendieck 主持；I.~Bucur、C.~Houzel、L.~Illusie、J.-P.~Jouanolou 及 J.-P.~Serre 协作。\par
\textbf{$\ell$ 进\Term{cohomology}{上同调}与 $L$ 函数}\par
Springer-Verlag，柏林、海德堡、纽约。

\SourcePage{2}
\textbf{原书内题名页}\par
《数学讲义》（Lecture Notes in Mathematics），A.~Dold 与 B.~Eckmann 主编，第 589 册。\par
布瓦玛丽\Term{algebraic-geometry}{代数几何}讨论班，1965—1966；SGA 5。\par
A.~Grothendieck 主持，I.~Bucur、C.~Houzel、L.~Illusie、J.-P.~Jouanolou 与 J.-P.~Serre 协作。\par
$\ell$ 进\Term{cohomology}{上同调}与 $L$ 函数。\par
Springer-Verlag，柏林、海德堡、纽约，1977 年。

\SourcePage{3}
\textbf{编者}\par
Luc Illusie\par
巴黎南大学（Université de Paris-Sud）\par
91405 Orsay，法国。\par
\medskip
AMS 学科分类（1970）：14C15、14F05、14F20、14G10、14G13、14H35、14M15、16A62、18F30。\par
ISBN 3-540-08248-4，Springer-Verlag，柏林、海德堡、纽约。\par
ISBN 0-387-08248-4，Springer-Verlag，纽约、海德堡、柏林。\par
\medskip
本书受版权保护。无论涉及全部材料还是其中一部分，一切权利均予保留，尤其包括翻译、重印、插图再利用、广播、使用复印机或类似手段复制，以及存入数据库的权利。

依照德国《著作权法》第 54 条，若所制作的复制件用于私人使用以外的目的，须向出版者支付费用，其数额由与出版者的协议确定。\par
\copyright{} Springer-Verlag，柏林、海德堡，1977 年。\par
德国印制。\par
印刷与装订：Beltz Offsetdruck，Hemsbach/Bergstr.\par
2141/3140-543210。

\SGAfrontchapter{引言}
\SourcePage{4}
本卷汇集了 Grothendieck 于 1965—1966 年在 IHÉS 主持的讨论班各讲；这些讲稿最初以油印笔记的\Term{form}{形式}分发。与初版相比，重要的变动只有两处：第 II 讲未收入本卷；第 III 讲则全部重写，并增添了一个编号为 III B 的附录。除少量细节上的修改及增补脚注外，其余各讲均保持原样。

本讨论班的核心是\Term{etale-cohomology}{平展上同调}中的 Lefschetz 公式（III、III B、XII），以及它在 $L$ 函数的\Term{cohomology}{上同调}解释中的应用（XIV）。Grothendieck 在 Bourbaki 讨论班报告 [2] 中宣布的全部结果，在这里都有完整证明。第 III、III B 和 XII 讲沿不同途径建立的迹公式，其一般性超过了仅为证明 $L$ 函数有理性所需的程度。Deligne 的报告（SGA $4\frac12$，Rapport）给出了这一有理性的一个明显更短而完整的证明；其中遵循 Grothendieck 在第 XII 与 XIV 讲中采用的方法，但舍去了所有不必要的一般性。不过，我们仍希望 III、III B 中的公式能够用于其他情形。至于本讨论班的其余部分，则包括两讲关于取\Term{limit}{极限}的理论——这一理论产生了\Term{ell-adic-cohomology}{$\ell$ 进上同调}（V、VI）——以及对\Term{duality-formalism}{对偶形式体系}（I、VII）和 Lefschetz 公式（VII、X）的若干补充。

本卷的内容，更确切地说，如下所述。

第 I 讲独立于讨论班的其余各讲。其主要结果是：在一个满足某些附加\Term{local-scheme}{局部}条件的\Term{regular-scheme}{正则概形}上，若能够使用\Term{resolution-of-singularities}{奇点消解}和\Term{purity-theorem}{纯性定理}，则当 $n$ 与所有\Term{residue-characteristic}{剩余特征}互素时，取值于 $\mathbb Z/n\mathbb Z$ 的\Term{constant-sheaf}{常值层}是\Term{dualizing-sheaf}{对偶化层}（I.3.4.1）。由 Hironaka 的工作和 Artin 的结果（SGA 4 XIX），这一\Term{theorem}{定理}尤其适用于\Term{characteristic-zero}{特征零}的\Term{excellent-scheme}{优良}\Term{regular-scheme}{正则概形}。

\SourcePage{5}
\Term{dimension}{维数}为 $1$ 的\Term{regular-scheme}{正则概形}情形另用一种非常初等的方法单独处理（也见 SGA $4\frac12$，Dualité）。还应指出，Deligne 在（SGA $4\frac12$，Th.~Finitude）中用另一种既不依赖\Term{resolution}{消解}也不依赖\Term{purity}{纯性}的论证证明：若 $f\colon X\longrightarrow S$ 是\Term{dimension}{维数}为 $0$ 或 $1$ 的\Term{regular-scheme}{正则概形}上的\Term{finite-type}{有限型}\Term{scheme}{概形}，则\Term{complex}{复形} $f^!(\mathbb Z/n\mathbb Z)_S$ 是\Term{dualizing-complex}{对偶化复形}（$n$ 与所有\Term{residue-characteristic}{剩余特征}互素）。这一结果比第 I 讲的\Term{biduality-theorem}{双对偶定理}更有用；它与后者有所交叠，但并不包含后者。不过，目前尚不知道它能否推广到基\Term{scheme}{概形} $S$ 为\Term{dimension}{维数}大于 $1$ 的\Term{excellent}{优良}\Term{regular-scheme}{正则概形}的情形。

如上所述，题为“任意\Term{support}{支撑}\Term{cohomology}{上同调}的 \Term{kunneth-formula}{Künneth 公式}”的第 II 讲未收入本卷。这一讲由 L.~Illusie 根据 Grothendieck 的手稿笔记整理，讨论的是泛\Term{cohomological-properness}{上同调固有性}与泛\Term{local-acyclicity}{局部无环性}\Term{theorem}{定理}，不过当时只在\Term{resolution}{消解}假设下给出了证明。Deligne 后来不借助这些假设证明了同样的结果（SGA $4\frac12$，Th.~Finitude），因此已无必要发表这一讲。某些补充，例如有关\Term{normal-crossings-divisor}{正规交叉除子}情形的补充，已收入该文的一个附录。此外，对于\Term{set-sheaf}{集合层}和非交换\Term{group-sheaf}{群层}，与该文结果类似的\Term{theorem}{定理}已由 Raynaud 夫人在她的报告（SGA 1 XIII）中证明。

第 III 讲证明了域上\Term{finite-type}{有限型}\Term{scheme}{概形}的层\Term{complex}{复形}之间的\Term{cohomological-correspondence}{上同调对应}所满足的 Lefschetz--Verdier 公式；第 III B 讲则说明如何在某些情形下计算公式中的\Term{local-term}{局部项}，例如\Term{curve}{曲线}上的 \Term{frobenius-correspondence}{Frobenius 对应}。关于这两讲内容的进一步说明，请读者参看各讲的引言。

A.~Grothendieck 的第 IV 讲讨论与一个\Term{cycle}{闭链}相关联的\Term{cohomology}{上同调}类的构造，已由 P.~Deligne 整理并发表于 SGA $4\frac12$。

\SourcePage{6}
第 V 讲先就\Term{adic-projective-system}{进投射系统}作技术准备，第 VI 讲随后定义\Term{constructible-ell-adic-sheaf}{可构造 $\mathbb Z_\ell$ 层}（以及\Term{constructible}{可构造}扭\Term{constant-sheaf}{常值层}；后者在 Deligne 的术语中称为“\Term{lisse-sheaf}{光滑}”层，见 [1]、SGA $4\frac12$，Arcata），并通过从 $\mathbb Z/\ell^n$ 层取\Term{limit}{极限}，将\Term{proper-support-higher-direct-image}{固有支撑高阶正像}的\Term{form}{形式}体系推广到\Term{constructible}{可构造} $\mathbb Z_\ell$ 层。然而，第 VI 讲没有定义 $\mathbb Z_\ell$ 层的\Term{derived-category}{导出范畴}，也没有在其中定义延拓 $\mathbb Z/\ell^n$ 层的\Term{derived-category}{导出范畴}上已有运算的 $\mathbin{\overset{\mathrm L}{\otimes}}$、$\mathrm R\underline{\operatorname{Hom}}$ 运算。因此，更没有考虑将 SGA 4 XVIII 的\Term{duality-formalism}{对偶形式体系}推广到 $\mathbb Z_\ell$ 层。这个问题是 Jouanolou 的博士论文（未发表）所研究的主题。不过，读者可参阅 [1] 的开头，那里定义了 $\mathbb Z_\ell$ 层的 $\underline{\operatorname{Ext}}^{\,i}$。

第 VII 讲独立于讨论班的其余部分，无疑也是其中最有意思的讲次之一。它包括：a）若干标准簇（\Term{affine-space}{仿射空间}、\Term{projective-space}{射影空间}、\Term{flag-variety}{旗簇}）的\Term{cohomology}{上同调}计算；b）\Term{etale-cohomology}{平展上同调}中\Term{chern-class}{陈类}理论的阐述；c）\Term{self-intersection-formula}{自交公式}
\[
 i^*i_*x=xc_d(N)
\]
在\Term{etale-cohomology}{平展上同调}中的证明（以及在\Term{chow-ring}{Chow 环}中的证明，后一情形的证明归于 Mumford），还有这一公式的各种应用（\Term{blow-up}{爆破}所得簇的\Term{cohomology}{上同调}计算、Gauss--Bonnet 公式）。

第 VII 讲包含关于\Term{derived-category}{导出范畴}中的\Term{finite}{有限}性条件以及\Term{grothendieck-group}{Grothendieck 群}的一系列一般论述。这些问题在（SGA 6 I、IV）中又在更一般的框架下重新讨论。
% 原书此处确印 VII；不擅改为 VIII。疑点登记于本块报告。

J.-P.~Serre 的第 IX 讲题为“Brauer 理论导论”，已发表于 [4]，本卷不再收入。其中一个主要结果（归于 Swan）涉及以“\Term{swan-character}{Swan 特征标}”为\Term{character}{特征标}的\Term{projective-module}{投射模}的存在性。这一结果用于第 X 讲，后者为\Term{curve}{曲线}上的层证明了一个 \Term{euler-poincare-formula}{Euler--Poincaré 公式}，比 Raynaud 的报告 [3] 中的公式稍为一般。

\SourcePage{7}
I.~Bucur 整理的第 XI 讲现已不存：它在一次搬家时遗失，而作者没有留副本。这一讲讨论的是 Grothendieck 的交换\Term{trace-theory}{迹理论}，它推广了 Stallings [5] 的理论。在（III B 5、6）中，以层论为框架，阐述了这一理论的一种更精细的\Term{form}{形式}。其在 Lefschetz 公式中的应用见第 XII 与 III B 讲。第 XII 讲的迹公式独立于第 III 讲的一般公式得到证明；不过，（III B 6）说明，其中的\Term{local-term}{局部项}确实就是一般公式中的\Term{local-term}{局部项}，而且这一一般公式蕴含第 XII 讲的公式。

第 XIII 讲的缺失纯粹缘于编号：第 X 至 XII 讲的内容，最初计划分四讲讲授。

第 XIV 讲（有时也编号为 XV）包含 \Term{frobenius-correspondence}{Frobenius 对应}的一般性质（特别是“几何”Frobenius 与“算术”Frobenius 的比较），并从第 XII 讲（或 III B）的迹公式出发，给出 Grothendieck 对 $L$ 函数有理性的证明。其中两个重要步骤是归约到\Term{curve}{曲线}情形，以及通过取\Term{limit}{极限}归约到关于 $\mathbb Z/\ell^n$ 层的一个公式。

在讨论班的口头报告中，Grothendieck 曾报告尚未解决的问题，并陈述了一些猜想。原定用来结束整个讨论班的这一讲，可惜没有写成讲稿；同样没有写成讲稿的，还有他那篇十分精彩的开场报告，其中回顾了各种背景（\Term{topology}{拓扑}、复解析、\Term{algebra}{代数}）下的 \Term{euler-poincare-formula}{Euler--Poincaré 公式}和 Lefschetz 公式。

感谢 P.~Deligne 说服我在第 III 讲的新稿中写出 Lefschetz--Verdier 公式的证明，从而消除了出版本讨论班讲稿的一项\Term{obstruction}{障碍}。我十分感谢他对行文提出的改进建议，

\SourcePage{8}
以及他在为出版者准备本卷时给予我的帮助。我也感谢 Lièvremont 夫人，她在很短的时间内细致地完成了第 III 与 III B 讲的打字工作。

在结束这篇引言之前，我还要向 I.~Bucur 致以追念；他于 1976 年 9 月因癌症去世。曾像我一样有幸与他为友的人，永远不会忘记他那无比的亲切、含笑的幽默，以及交谈中流露的魅力。

\begin{flushright}
巴黎，1977 年 2 月 19 日\par
Luc Illusie
\end{flushright}

\SGAfrontchapter{引言参考文献}
\SourcePage{9}
\begin{enumerate}
\item[{[1]}] Deligne, P.，\emph{Weil 猜想 II}（La conjecture de Weil II），待刊。
\item[{[2]}] Grothendieck, A.，\emph{Lefschetz 公式与 $L$ 函数的有理性}（Formule de Lefschetz et rationalité des fonctions $L$），Bourbaki 讨论班，1964—1965 年，第 279 讲；收入\emph{关于\Term{scheme}{概形}\Term{cohomology}{上同调}的十篇报告}（Dix exposés sur la cohomologie des schémas），North Holland Pub. Co.，1968 年。
\item[{[3]}] Raynaud, M.，\emph{层的 \Term{euler-poincare-characteristic}{Euler--Poincaré 特征}与\Term{abelian-variety}{阿贝尔簇}的\Term{cohomology}{上同调}（根据 Ogg--Shafarevitch 与 Grothendieck 的工作）}（Caractéristique d'Euler-Poincaré d'un faisceau et cohomologie des variétés abéliennes (d'après Ogg-Shafarevitch et Grothendieck)），Bourbaki 讨论班，1964—1965 年，第 286 讲；收入\emph{关于\Term{scheme}{概形}\Term{cohomology}{上同调}的十篇报告}，North Holland Pub. Co.，1968 年。
\item[{[4]}] Serre, J.-P.，\emph{\Term{finite-group}{有限群}的\Term{linear-representation}{线性表示}}（Représentations linéaires des groupes finis），Hermann，修订第 2 版，1971 年。
\item[{[5]}] Stallings, J.，\emph{无\Term{center}{中心}群——Gottlieb \Term{theorem}{定理}的一个\Term{algebra}{代数}表述}（Centerless groups---an algebraic formulation of Gottlieb's theorem），\emph{Topology}，第 4 卷，第 129--134 页，1965 年。
\end{enumerate}

\SGAfrontchapter{SGA 书目}
\SourcePage{10}
SGA 是“布瓦玛丽\Term{algebraic-geometry}{代数几何}讨论班”（Séminaire de Géométrie Algébrique du Bois-Marie）的缩写。
\begin{description}
\item[SGA 1：] \emph{\Term{etale-covering}{平展覆盖}与\Term{fundamental-group}{基本群}}（Revêtements étales et groupe fondamental），A.~Grothendieck；Lecture Notes in Mathematics，第 224 册，Springer-Verlag，1971 年。
\item[SGA 2：] \emph{\Term{coherent-sheaf}{凝聚层}的\Term{local-scheme}{局部}\Term{cohomology}{上同调}与\Term{local-scheme}{局部}及整体 Lefschetz \Term{theorem}{定理}}（Cohomologie locale des faisceaux cohérents et théorèmes de Lefschetz locaux et globaux），A.~Grothendieck；North Holland Pub. Co.，1968 年。
\item[SGA 3：] \emph{\Term{group-scheme}{群概形} I、II、III}（Schémas en groupes I, II, III），M.~Demazure 与 A.~Grothendieck；Lecture Notes in Mathematics，第 151、152、153 册，Springer-Verlag，1970 年。
\item[SGA 4：] \emph{\Term{topos}{拓扑斯}理论与\Term{scheme}{概形}的\Term{etale-cohomology}{平展上同调}}（Théorie des topos et cohomologie étale des schémas），M.~Artin、A.~Grothendieck、J.-L.~Verdier；Lecture Notes in Mathematics，第 269、270、305 册，Springer-Verlag，1972—1973 年。
\item[SGA $4\frac12$：] \emph{\Term{etale-cohomology}{平展上同调}}（Cohomologie étale），P.~Deligne；Lecture Notes in Mathematics，第 569 册，Springer-Verlag，1977 年。
\item[SGA 5：] \emph{$\ell$ 进\Term{cohomology}{上同调}与 $L$ 函数}（Cohomologie $\ell$-adique et fonctions $L$），A.~Grothendieck；Lecture Notes in Mathematics，第 589 册，Springer-Verlag，1977 年。
\item[SGA 6：] \emph{\Term{intersection-theory}{相交理论}与 \Term{riemann-roch-theorem}{Riemann--Roch 定理}}（Théorie des intersections et théorème de Riemann-Roch），P.~Berthelot、A.~Grothendieck、L.~Illusie；Lecture Notes in Mathematics，第 225 册，Springer-Verlag，1971 年。
\item[SGA 7：] \emph{\Term{algebraic-geometry}{代数几何}中的\Term{monodromy-group}{单值群}}（Groupes de monodromie en géométrie algébrique），第一部分 A.~Grothendieck，第二部分 P.~Deligne 与 N.~Katz；Lecture Notes in Mathematics，第 288、340 册，Springer-Verlag，分别于 1972、1973 年出版。
\end{description}

\SGAfrontchapter{原书目录标题}
\SourcePage{11}
\noindent 目录\footnote{关于第 II、IV、IX、XI、XIII 讲缺失的说明，见引言。}\par\medskip
% 下列各项仅保存原目录的标题与署名；译本页码由总工程目录重建。
\noindent\textbf{第 I 讲}\quad \Term{dualizing-complex}{对偶化复形}。A.~Grothendieck 主讲，L.~Illusie 整理。\par\medskip
\noindent\textbf{第 III 讲}\quad Lefschetz 公式。A.~Grothendieck 主讲，L.~Illusie 整理。\par\medskip
\noindent\textbf{第 III B 讲}\quad \Term{local-term}{局部项}的计算。L.~Illusie。\par\medskip
\noindent\textbf{第 V 讲}\quad $J$ \Term{adic-projective-system}{进投射系统}。J.-P.~Jouanolou。\par\medskip
\noindent\textbf{第 VI 讲}\quad $\ell$ 进\Term{cohomology}{上同调}。J.-P.~Jouanolou。\par\medskip
\noindent\textbf{第 VII 讲}\quad 若干经典\Term{scheme}{概形}的\Term{cohomology}{上同调}与\Term{chern-class}{陈类}的\Term{cohomology}{上同调}理论。J.-P.~Jouanolou。\par\medskip
\noindent\textbf{第 VIII 讲}\quad \Term{abelian-category}{阿贝尔范畴}与三角\Term{category}{范畴}的类群，\Term{perfect}{完美}\Term{complex}{复形}。A.~Grothendieck 主讲，I.~Bucur 整理。\par\medskip
\noindent\textbf{第 X 讲}\quad \Term{etale-cohomology}{平展上同调}中的 \Term{euler-poincare-formula}{Euler--Poincaré 公式}。A.~Grothendieck 主讲，I.~Bucur 整理。\par\medskip

\SourcePage{12}
\noindent\textbf{第 XII 讲}\quad \Term{algebraic-geometry}{代数几何}中的 Nielsen--Wecken 公式与 Lefschetz 公式。A.~Grothendieck 主讲，I.~Bucur 整理。\par\medskip
\noindent\textbf{第 XIV（即 XV）讲}\quad Frobenius \Term{morphism}{态射}与 $L$ 函数的有理性。C.~Houzel。\par\medskip
\noindent\textbf{术语索引}\par\medskip
\noindent\textbf{记号索引}\par\medskip

\begin{center}
\begin{tikzpicture}[>=stealth, every node/.style={font=\small}]
\node (left) at (0,0) {I，VII，VIII};
\node (x) at (3,0) {X};
\draw[->] (left.east) -- (x.west);
\node at (3.6,0) {，};
\node (iii) at (4.5,0.8) {III};
\node (iiib) at (6.7,0.8) {III B};
\node (v) at (4.5,0) {V};
\node (vi) at (6.7,0) {VI};
\node (iiibfive) at (4.5,-0.8) {III B 5};
\node (xii) at (6.7,-0.8) {XII};
\node (xiv) at (9,0) {XIV};
\draw[->] (iii) -- (iiib);
\draw[->] (v) -- (vi);
\draw[->] (iiibfive) -- (xii);
\draw[->] (iiib.east) -- (xiv.north west);
\draw[->] (vi) -- (xiv);
\draw[->] (xii.east) -- (xiv.south west);
\end{tikzpicture}
\end{center}
